\documentclass{article}

\PassOptionsToPackage{numbers,sort&compress}{natbib}
\usepackage[preprint]{neurips_2026}

\usepackage[utf8]{inputenc}
\usepackage[T1]{fontenc}
\usepackage{hyperref}
\usepackage{url}
\usepackage{booktabs}
\usepackage{amsmath,amssymb,amsfonts,amsthm}
\usepackage{enumitem}
\usepackage{microtype}
\setlength{\emergencystretch}{2em}

\theoremstyle{definition}
\newtheorem{definition}{Definition}[section]
\newtheorem{problem}[definition]{Problem}
\newtheorem{assumption}[definition]{Assumption}
\theoremstyle{plain}
\newtheorem{lemma}[definition]{Lemma}
\newtheorem{theorem}[definition]{Theorem}
\newtheorem{proposition}[definition]{Proposition}
\newtheorem{corollary}[definition]{Corollary}
\theoremstyle{remark}
\newtheorem{remark}[definition]{Remark}

\DeclareMathOperator*{\argmin}{arg\,min}
\DeclareMathOperator{\KL}{KL}
\newcommand{\E}{\mathbb{E}}
\newcommand{\Prob}{\mathbb{P}}

\title{How to Pick a Model}
\author{Anonymous Authors}

\begin{document}

\maketitle

\section{Introduction}
\label{sec:introduction}

Which model should be deployed on a new task whose solution can be checked?
Current practice often answers this question indirectly: evaluate candidate
systems end to end on a benchmark and deploy the one with the best aggregate
score.  We call this practice \emph{benchmark racing}.  It has produced
increasingly realistic evaluations, from
repository repair against executable tests to machine-learning engineering
against held-out validation criteria~\citep{jimenez2024swebench,chan2024mlebench}.
Yet an aggregate score is not sufficient for deployment when models
differ in serving cost, can be retried, or exhibit complementary strengths
across instances.  A system that is more accurate per attempt need not be the
one that reaches the first verified solution with the least wall-clock time,
computation, energy, tokens, or money.  Moreover, the same aggregate advantage
can arise from cheaper attempts, greater competence on the relevant kind of
instance, more informative pre-deployment evidence, or better allocation of
effort.  These mechanisms imply different engineering decisions, but an
end-to-end race reports only their net result.

The problem has important antecedents.  Classical algorithm selection maps
instance features to a solver from a portfolio; empirical performance models
predict runtimes on unseen instances; and systems such as SATzilla demonstrate
that the best solver can be instance dependent rather than
global~\citep{rice1976algorithm,hutter2014runtime,xu2012satzilla,kotthoff2016algorithm}.
Statistical racing and best-arm methods reduce the cost of identifying a winner
from repeated evaluations~\citep{maron1993hoeffding,maron1997racing,kaufmann2016bestarm}.
More recently, LLM cascades and routers have learned cost--quality policies over
heterogeneous model pools~\citep{chen2023frugalgpt,ong2024routellm,hu2024routerbench},
while test-time scaling and sample-allocation studies show that model size,
sampling policy, and inference effort interact in task- and budget-dependent
ways~\citep{brown2024monkeys,snell2025testtime,wu2025inference,osca2024}.
These lines of work establish the relevance of per-instance choice and
resource-aware inference.  Their primary output is a runtime predictor,
selector, or allocation policy under a declared performance objective.  Our
focus is complementary: an accounting law for the verified first-passage
resource of a stochastic solver that separates the price of one attempt, its
mode-conditional probability of success, the information available before
deployment, and the extent to which the scheduler uses that information.

Our conceptual foundation is the proper-time framework of Achille and
Soatto~\citep{achille2025agents}.  They model an agent as a stochastic dynamical
system acting on verifiable transductive tasks and argue that, because a
universal search procedure may eventually find a valid witness, accuracy alone
does not characterize useful intelligence.  The relevant benefit of learning
is instead to reduce the time required to find a solution.  Their proper time
for reaching a target state is a path quantity,
\(
  \tau_\nu(u\!\to\!v)=\min_{h:u\to v} T(h)/\nu(h\mid u)
\),
which divides trajectory length by the probability of realizing that path.  We
retain two principles from this theory: solutions must be
defined by an external verifier, and information is valuable insofar as it buys
speed.  We do not identify their path proper time with our deployment metric.
Their framework supplies a computational foundation rather than a finite-sample
rule for selecting from a deployed model menu.  For a design~\(d\), instance
\(X\sim\mu\), and run randomness~\(\xi\), we define certified
resource-to-solution \(T_d(\delta)\) as the \((1-\delta)\)-quantile of the
resource consumed before the first verified artifact.  This is a measurable
service-level objective for comparing a finite menu under one predeclared
resource clock.

For a tractable accounting, we introduce a solver-relevant mode \(S\),
admissible evidence \(Z\), and a physically implemented effort share \(q_Z(S)\).
The packed hypothesis makes mode-conditioned resource proportional to unit
attempt cost \(\kappa(M,S)\), focused scale \(t_0(M,S)\), and \(1/q_Z(S)\).
Under a prior-matched baseline, the resulting expected log-resource gain is
\begin{equation}
\begin{aligned}
\Delta_{\mathrm{pack}}
={}&
\underbrace{\mathbb E_S\!\left[\log
  \frac{\kappa(M_0,S)}{\kappa(M,S)}\right]}_{\text{unit cost}}
+
\underbrace{\mathbb E_S\!\left[\log
  \frac{t_0(M_0,S)}{t_0(M,S)}\right]}_{\text{focused competence}}
\\[-0.15em]
&+
\underbrace{I(S;Z)}_{\text{mode information}}
-
\underbrace{\mathbb E_Z\!\left[
  \mathrm{KL}(\pi_Z\Vert q_Z)\right]}_{\text{allocation mismatch}}.
\label{eq:intro-four-term}
\end{aligned}
\end{equation}
The identity is exact inside the packed model; it is not a general
decomposition of \(\log T_d(\delta)\).  We therefore derive the corresponding
residual form when inverse-share log-linearity is approximate and an exact
geometric-retry bridge showing that, under IID eligible attempts, \(1/q\)
governs mean first-passage attempts, whereas certified quantiles require an
explicit correction.  This distinction makes the assumption testable rather
than definitional.

Equation~\eqref{eq:intro-four-term} turns model selection into a diagnostic and
counterfactual calculation: it separates cheaper attempts, model competence,
information acquisition, and allocation quality.  A calibrated pilot can then
predict held-out expected log first-passage resource, while full deployment
trials remain the confirmation standard.

We make three contributions:
\begin{enumerate}[leftmargin=1.5em,itemsep=0.2em,topsep=0.35em]
  \item Starting from Achille and Soatto's time-centric principle, we formulate
  finite, verifiable model selection through first-passage resource and its
  certified population quantile, while keeping this deployment object distinct
  from both
  Achille--Soatto path proper time and the theorem-facing expected-log
  functional.
  \item We prove the four-term cost--competence--information--mismatch identity,
  its approximate-packed residual guarantee, and model-selection and retry
  consequences.  Together these results specify which quantities must be
  estimated and the assumptions under which their sum predicts a packed
  expected-log deployment comparison.
  \item We test those assumptions with real language models on a controlled,
  execution-verified coding benchmark.  In a held-out study of 175,104 physical
  trajectories over 768 unseen tasks using Qwen2.5-Coder 7B and 14B, the
  inverse-share exponent is \(0.987\), the four-term prediction has residual RMS
  \(0.015\) nat, and its model choice agrees with the held-out resource oracle in
  all six conditions.  Because a strict residual omnibus rejects exact closure,
  we interpret the decomposition as a falsifiable decision approximation, not
  a universal equality.
\end{enumerate}

\section{Verifier-Backed Model Selection}
\label{sec:problem}

We first separate the deployment quantity that the builder ultimately cares
about from the expected-log functional that admits the decomposition.  This
distinction fixes the task, the resource unit, and the statistical estimand
before introducing latent modes.

\begin{definition}[Verifiable task family and resource clock]
\label{def:task-family}
A \emph{verifiable task family} is a tuple
\(
  \mathcal T=(\mathcal X,\mathcal Y,V,\mu)
\), where \(X\sim\mu\) is a deployment instance, \(y\in\mathcal Y\) is a
candidate artifact, and \(V:\mathcal X\times\mathcal Y\to\{0,1\}\) is a
predeclared external verifier.  Before comparing systems, the builder fixes one
scalar, additive resource clock---for example isolated GPU-seconds, joules,
tokens, or dollars.  Signal acquisition, generation, verification, and retry
costs are charged whenever they consume that resource.
\end{definition}

A design \(d\) fixes the model, serving configuration, admissible signal, and
allocation or retry policy.  Its randomness is denoted by \(\xi\), and
\(Y_d(t;x,\xi)\) is the latest artifact available by resource time \(t\).

\begin{definition}[First passage and certified resource-to-solution]
\label{def:certified-time}
The first verified passage of design \(d\) is
\begin{equation}
  \tau_d(x,\xi)
  :=\inf\{t\ge0:V(x,Y_d(t;x,\xi))=1\},
  \qquad \inf\varnothing:=\infty.
  \label{eq:first-passage}
\end{equation}
At failure tolerance \(\delta\in(0,1)\), its certified
resource-to-solution is the population quantile
\begin{equation}
  T_d(\delta)
  :=\inf\{t\ge0:
  \Prob_{X\sim\mu,\xi}[\tau_d(X,\xi)\le t]\ge1-\delta\}.
  \label{eq:certified-time}
\end{equation}
For a baseline \(d_0\), the certified log-speedup is
\(\Delta_{\mathrm{cert}}(d_0,d;\delta)
=\log T_{d_0}(\delta)-\log T_d(\delta)\), whenever both quantiles are finite
and positive.
\end{definition}

\begin{problem}[Deployment model selection]
\label{prob:model-selection}
Given a finite, predeclared design menu \(\mathcal D\), a task family, a common
resource clock, and \(\delta\), select
\begin{equation}
  d^\star\in\argmin_{d\in\mathcal D}T_d(\delta).
  \label{eq:deployment-selection}
\end{equation}
The comparison is undefined if systems use different clocks or verifiers; a
multi-resource problem must instead declare a scalarization or a Pareto
frontier before selection.
\end{problem}

This formulation covers code against tests, proof search against a checker, and
other transductive tasks in which one concrete artifact must be produced for
the observed instance.  It does not assume that a successful artifact is
unique.

\subsection{Latent modes and packed allocation}
\label{sec:packed-setup}

\begin{definition}[Modes, signals, and deployable shares]
\label{def:packed-objects}
Let \(S\in\mathcal S=\{1,\ldots,m\}\) be a solver-relevant latent mode with
prior \(\pi\).  Mode semantics are fixed independently of the model being
compared.  An admissible signal \(Z\) is observed before the allocated work is
executed, with posterior \(\pi_z(s):=\Prob[S=s\mid Z=z]\).  A \emph{packed
allocation} is a measurable map \(z\mapsto q_z\in\Delta(\mathcal S)\), where
\(q_z(s)\) is the physically implemented share of eligible attempts devoted to
mode-\(s\) work.  We require \(q_z(s)>0\) whenever \(\pi_z(s)>0\).

For model \(M\) and mode \(s\), \(\kappa(M,s)>0\) is the expected resource of
one eligible attempt under the frozen serving configuration, and
\(t_0(M,s)>0\) is its dimensionless focused first-passage scale when all
eligible effort is assigned to that mode.
\end{definition}

For distributions \(a,b\in\Delta(\mathcal S)\), write
\[
  H(a):=-\sum_s a(s)\log a(s),
  \qquad
  \KL(a\Vert b):=\sum_s a(s)\log\frac{a(s)}{b(s)},
\]
with \(0\log0:=0\), and define
\(I(S;Z):=H(\pi)-\E_ZH(\pi_Z)\).

The signal may be a static instance feature or a charged partial trace.  The
allocation is a resource-sharing mechanism, not merely a label predicted by a
classifier.

\begin{assumption}[Packed inverse-share law]
\label{ass:packed}
For a shared constant \(c_{\mathrm{sim}}>0\), the mode-conditioned resource
scale of \((M,q)\) is
\begin{equation}
  \mathsf T_{M,q}(s,z)
  =c_{\mathrm{sim}}\,
    \kappa(M,s)\frac{t_0(M,s)}{q_z(s)}.
  \label{eq:packed-law}
\end{equation}
The law applies only when the scheduler implements the declared shares.  Any
model- or mode-dependent overhead not represented by \(\kappa\) violates the
exact assumption and must enter the residual model below.
\end{assumption}

In particular, a one-time signal-acquisition cost is part of the actual
resource clock.  Exact packedness requires that this charge already scale with
the multiplicative product in~\eqref{eq:packed-law}; an additive or
allocation-dependent charge cannot generally be absorbed into \(\kappa(M,s)\)
without changing the factorization and therefore appears in the residual of
Proposition~\ref{prop:approx-packed}.

\begin{definition}[Packed expected-log risk]
\label{def:packed-risk}
For the joint law of \((S,Z)\), define
\begin{equation}
  \mathcal L(M,q):=\E[\log\mathsf T_{M,q}(S,Z)].
  \label{eq:packed-risk}
\end{equation}
For a baseline model \(M_0\) with fixed allocation \(q_0\), define
\(
  \Delta_{\mathrm{pack}}
  :=\mathcal L(M_0,q_0)-\mathcal L(M,q)
\).
We require \(q_0(s)>0\) whenever \(\pi(s)>0\).
\end{definition}

\begin{remark}[Metric and identifiability guardrails]
\label{rem:guardrails}
In general,
\(
  \Delta_{\mathrm{pack}}
  \ne\Delta_{\mathrm{cert}}(d_0,d;\delta)
\), and \(\E[\log T]\ne\log\E[T]\).  Section~\ref{sec:quantile-bridge}
gives an exact bridge in an IID geometric setting; outside such a bridge, the
two quantities must be reported separately.  Moreover, total resource
identifies only the product \(\kappa(M,s)t_0(M,s)\).  Interpreting its factors
as unit cost and competence requires separate attempt-level measurement of
\(\kappa\).
\end{remark}

\section{Four-Term Speed Accounting}
\label{sec:accounting}

The decomposition follows from the conditional cross-entropy incurred by the
physical allocation.  Complete proofs of every result in this section and
Section~\ref{sec:quantile-bridge} appear in Appendix~\ref{app:proofs}.

\begin{lemma}[Conditional log-resource and posterior matching]
\label{lem:conditional}
Under Assumption~\ref{ass:packed}, for any signal value \(z\),
\begin{equation}
\begin{aligned}
\E[\log\mathsf T_{M,q}(S,z)\mid Z=z]
={}&\E[\log(c_{\mathrm{sim}}\kappa(M,S)t_0(M,S))\mid Z=z]
\\
&+H(\pi_z)+\KL(\pi_z\Vert q_z).
\end{aligned}
\label{eq:conditional-decomposition}
\end{equation}
Consequently, \(q_z=\pi_z\) uniquely minimizes conditional expected
log-resource, and its excess over the minimum is exactly
\(\KL(\pi_z\Vert q_z)\).
\end{lemma}

\begin{theorem}[Model-aware four-term decomposition]
\label{thm:four-term}
Let \((M_0,q_0)\) use a fixed baseline allocation and let \((M,q)\) use signal
\(Z\).  Under Assumption~\ref{ass:packed}, define
\begin{align*}
  C(M_0,M)&:=\E_S\!\left[\log\frac{\kappa(M_0,S)}{\kappa(M,S)}\right],
  &
  \Phi(M_0,M)&:=\E_S\!\left[\log\frac{t_0(M_0,S)}{t_0(M,S)}\right],\\
  G(M)&:=I(S;Z),
  &
  \varepsilon(M,q)&:=\E_Z[\KL(\pi_Z\Vert q_Z)].
\end{align*}
Then
\begin{equation}
  \boxed{
  \Delta_{\mathrm{pack}}
  =C(M_0,M)+\Phi(M_0,M)+G(M)
   +\KL(\pi\Vert q_0)-\varepsilon(M,q).}
  \label{eq:four-term-general}
\end{equation}
For the prior-matched baseline \(q_0=\pi\), the baseline mismatch vanishes,
leaving the four terms \emph{unit cost}, \emph{focused competence},
\emph{mode information}, and \emph{allocation mismatch}.
\end{theorem}

The identity is exact because Assumption~\ref{ass:packed} makes log-resource
additive in these quantities.  Its value is diagnostic: it turns an aggregate
gain into separately measurable mechanisms and exposes the assumption that
must survive held-out residual tests.

\begin{corollary}[Model-selection score]
\label{cor:model-score}
For candidate systems sharing \((\pi,\mathcal S,c_{\mathrm{sim}})\), possibly
with candidate-specific signals \(Z_M\), posteriors
\(\pi_z^M(s):=\Prob[S=s\mid Z_M=z]\), and allocations \(q^M\), define
\begin{equation}
\begin{aligned}
  \mathcal J(M,q^M)
  :={}&\E_S[\log\kappa(M,S)]
      +\E_S[\log t_0(M,S)]
      -I(S;Z_M)\\
      &+\E_{Z_M}[\KL(\pi^{M}_{Z_M}\Vert q^M_{Z_M})].
\end{aligned}
\label{eq:model-score}
\end{equation}
Then \(\mathcal L(M,q^M)=\log c_{\mathrm{sim}}+H(\pi)+\mathcal J(M,q^M)\).
Thus the packed-optimal candidate is any minimizer of \(\mathcal J\).
\end{corollary}

Exact packedness is not needed for a useful decision, but its failure must be
quantified rather than hidden.

\begin{proposition}[Approximate packedness]
\label{prop:approx-packed}
Suppose the actual mode-conditioned scale satisfies
\begin{equation}
  \log\widetilde{\mathsf T}_{M,q}(s,z)
  =\log\mathsf T_{M,q}(s,z)+r_M(s,z,q),
  \qquad |r_M(s,z,q)|\le\eta
  \label{eq:approx-packed}
\end{equation}
uniformly for both compared systems.  For a prior-matched baseline,
\begin{equation}
  \widetilde\Delta
  =\underbrace{C+\Phi+G-\varepsilon}_{\Delta_4}+R,
  \qquad
  R=\E[r_{M_0}(S,q_0)-r_M(S,Z,q)],
  \qquad |R|\le2\eta.
  \label{eq:residual-decomposition}
\end{equation}
Here \(r_{M_0}(s,q_0)\) denotes the baseline residual, which has no signal
argument because \(q_0\) is fixed.
\end{proposition}

\begin{corollary}[Residual-aware decision margin]
\label{cor:robust-margin}
Let \(\widehat\Delta_4\) estimate the four-term prediction with
\(|\widehat\Delta_4-\Delta_4|\le e\), and suppose \(|R|\le\rho\).
Then \(\widetilde\Delta>0\) whenever
\(\widehat\Delta_4>e+\rho\), and \(\widetilde\Delta<0\) whenever
\(\widehat\Delta_4<-e-\rho\).  Inside the band
\(|\widehat\Delta_4|\le e+\rho\), the accounting does not certify a winner.
\end{corollary}

This margin is the decision-theoretic output of the theory: estimation error
and structural residual are both charged before changing the deployed model.

\section{From First-Passage Means to Certified Quantiles}
\label{sec:quantile-bridge}

The inverse-share factor \(1/q\) has an exact first-passage interpretation, but
the interpretation depends on whether the target is a mean or a quantile.

\begin{theorem}[Geometric allocation bridge]
\label{thm:geometric-bridge}
Fix one mode.  In each global slot, independently allocate an eligible attempt
to the true mode with probability \(q\in(0,1]\).  A true-mode attempt succeeds
independently with probability \(p\in(0,1)\), and off-mode attempts cannot
succeed.  If \(N_{p,q}\) is the first successful global slot, then
\begin{equation}
  N_{p,q}\sim\operatorname{Geom}(pq),
  \qquad
  \E[N_{p,q}]=\frac{1}{pq}.
  \label{eq:geometric-mean}
\end{equation}
Its certified slot quantile is
\begin{equation}
  Q_{p,q}(\delta)
  =\left\lceil\frac{\log\delta}{\log(1-pq)}\right\rceil.
  \label{eq:geometric-quantile}
\end{equation}
Let \(h(x):=-\log(1-x)\) and remove the integer ceiling, so
\(\widetilde Q_{p,q}(\delta)=-\log\delta/h(pq)\).  Comparing this quantity with
the inverse-share extrapolation of its focused value,
\(\widetilde Q_{p,1}(\delta)/q\), gives the log residual
\begin{equation}
  r_{\mathrm{inv}}(p,q)
  :=\log\frac{\widetilde Q_{p,q}(\delta)}
                    {\widetilde Q_{p,1}(\delta)/q}
  =\log\frac{q\,h(p)}{h(pq)},
  \qquad
  0\le r_{\mathrm{inv}}(p,q)\le-\log(1-p).
  \label{eq:quantile-residual}
\end{equation}
For the unrounded quantile \(\widetilde Q_{p,q}\), ceiling contributes at most
\begin{equation}
  0\le\log Q_{p,q}(\delta)-\log\widetilde Q_{p,q}(\delta)
  <\log\!\left(1+\frac{1}{\widetilde Q_{p,q}(\delta)}\right).
  \label{eq:ceiling-residual}
\end{equation}
\end{theorem}

Thus \(1/q\) is exact for mean first passage and a controlled approximation for
a cellwise quantile whose log residual vanishes as \(p\to0\).  For heterogeneous
modes or success
probabilities, the population quantity in Definition~\ref{def:certified-time}
is a mixture quantile; it cannot in general be obtained by averaging cellwise
quantiles.  Accordingly, \(t_0^{\mathrm{mean}}=1/p\) and the focused quantile
in~\eqref{eq:geometric-quantile} are different estimands and must not be
interchanged.

\begin{corollary}[Single-mode retry crossover]
\label{cor:retry-crossover}
Let a high-cost system have one-attempt success \(p_h\) and cost \(\kappa_h\),
and a lower-cost system have \(0<p_\ell<p_h<1\) and per-attempt cost
\(\kappa_\ell\).  The minimum retry depth whose fixed block matches or exceeds
the high-cost system's reliability is
\begin{equation}
  D_\times
  =\left\lceil\frac{\log(1-p_h)}{\log(1-p_\ell)}\right\rceil.
  \label{eq:retry-crossover}
\end{equation}
The lower-cost block is cheaper when \(D_\times\kappa_\ell<\kappa_h\).  If it
stops at first success, its expected cost through depth \(D\) is instead
\begin{equation}
  \kappa_\ell\frac{1-(1-p_\ell)^D}{p_\ell}.
  \label{eq:stopped-retry-cost}
\end{equation}
\end{corollary}

Equations~\eqref{eq:four-term-general}--\eqref{eq:stopped-retry-cost} provide
the complete theoretical chain used by the paper: define the deployment
target, decompose the expected-log packed surrogate, bound structural error,
and certify a decision only outside the combined uncertainty margin.
\label{sec:theory-end}

\bibliographystyle{plainnat}
\bibliography{references}

\clearpage
\appendix

\section{Proofs}
\label{app:proofs}

\subsection{Proof of Lemma~\ref{lem:conditional}}

\begin{proof}
Taking logs in~\eqref{eq:packed-law} and conditioning on \(Z=z\) gives
\begin{align*}
\E[\log\mathsf T_{M,q}(S,z)\mid Z=z]
={}&\E[\log(c_{\mathrm{sim}}\kappa(M,S)t_0(M,S))\mid Z=z]\\
&+\sum_{s\in\mathcal S}\pi_z(s)(-\log q_z(s)).
\end{align*}
The final sum is the cross-entropy between \(\pi_z\) and \(q_z\), and
\begin{align*}
\sum_s\pi_z(s)(-\log q_z(s))
&=-\sum_s\pi_z(s)\log\pi_z(s)
  +\sum_s\pi_z(s)\log\frac{\pi_z(s)}{q_z(s)}\\
&=H(\pi_z)+\KL(\pi_z\Vert q_z).
\end{align*}
For completeness, Jensen's inequality for \(-\log\) gives
\[
\KL(\pi_z\Vert q_z)
=\E_{\pi_z}\!\left[-\log\frac{q_z(S)}{\pi_z(S)}\right]
\ge-\log\!\left(\sum_{s:\pi_z(s)>0}q_z(s)\right)\ge0.
\]
Equality requires all mass to lie on the posterior support and
\(q_z(s)/\pi_z(s)\) to be constant there, hence \(q_z=\pi_z\).  The support
condition in Definition~\ref{def:packed-objects} makes every displayed quantity
finite.
\end{proof}

\subsection{Proof of Theorem~\ref{thm:four-term}}

\begin{proof}
For the deployed system, Lemma~\ref{lem:conditional} and the tower property
give
\begin{align}
\mathcal L(M,q)
={}&\log c_{\mathrm{sim}}
 +\E_S[\log\kappa(M,S)]
 +\E_S[\log t_0(M,S)]\nonumber\\
&+\E_Z[H(\pi_Z)]+\varepsilon(M,q).
\label{eq:proof-deployed-risk}
\end{align}
The fixed baseline allocation has cross-entropy
\begin{equation}
  \E_S[-\log q_0(S)]=H(\pi)+\KL(\pi\Vert q_0),
  \label{eq:proof-baseline-crossentropy}
\end{equation}
and hence
\begin{align}
\mathcal L(M_0,q_0)
={}&\log c_{\mathrm{sim}}
 +\E_S[\log\kappa(M_0,S)]
 +\E_S[\log t_0(M_0,S)]\nonumber\\
&+H(\pi)+\KL(\pi\Vert q_0).
\label{eq:proof-baseline-risk}
\end{align}
Subtracting~\eqref{eq:proof-deployed-risk} from
\eqref{eq:proof-baseline-risk}, and using
\(
I(S;Z)=H(\pi)-\E_ZH(\pi_Z)
\), yields~\eqref{eq:four-term-general}.  If \(q_0=\pi\), then
\(\KL(\pi\Vert q_0)=0\).
\end{proof}

\subsection{Proof of Corollary~\ref{cor:model-score}}

\begin{proof}
Apply~\eqref{eq:proof-deployed-risk} to each candidate and substitute
\(
\E_{Z_M}H(\pi^M_{Z_M})=H(\pi)-I(S;Z_M)
\).  The terms \(\log c_{\mathrm{sim}}+H(\pi)\) are common to every candidate;
the remaining terms are exactly \(\mathcal J(M,q^M)\).  Therefore the
minimizers of \(\mathcal L\) and \(\mathcal J\) coincide.
\end{proof}

\subsection{Proof of Proposition~\ref{prop:approx-packed}}

\begin{proof}
Substitute~\eqref{eq:approx-packed} into the expected log-resource difference.
The ideal terms equal \(\Delta_4\) by Theorem~\ref{thm:four-term}; the remaining
terms are
\(
R=\E[r_{M_0}(S,q_0)-r_M(S,Z,q)]
\).  The triangle inequality gives
\[
  |R|
  \le\E|r_{M_0}(S,q_0)|+\E|r_M(S,Z,q)|
  \le2\eta.
\]
\end{proof}

\subsection{Proof of Corollary~\ref{cor:robust-margin}}

\begin{proof}
Write \(a=\widehat\Delta_4-\Delta_4\).  Then
\(
\widetilde\Delta=\widehat\Delta_4-a+R
\), where \(|a|\le e\) and \(|R|\le\rho\).  Hence
\[
  \widehat\Delta_4-e-\rho
  \le\widetilde\Delta
  \le\widehat\Delta_4+e+\rho.
\]
The two sign claims follow from the corresponding strict endpoint inequality.
If the interval contains zero, these bounds alone cannot determine the sign.
\end{proof}

\subsection{Proof of Theorem~\ref{thm:geometric-bridge}}

\begin{proof}
A global slot succeeds exactly when it is allocated to the true mode and that
eligible attempt succeeds.  Independence therefore gives per-slot success
probability \(pq\), so
\[
  \Prob[N_{p,q}=n]=(1-pq)^{n-1}pq,
  \qquad
  \Prob[N_{p,q}>n]=(1-pq)^n.
\]
The survival-sum formula yields
\[
  \E[N_{p,q}]
  =\sum_{n\ge0}\Prob[N_{p,q}>n]
  =\sum_{n\ge0}(1-pq)^n
  =\frac{1}{pq}.
\]
Moreover,
\(
\Prob[N_{p,q}\le n]=1-(1-pq)^n
\).  Requiring this probability to be at least \(1-\delta\) gives
\(
(1-pq)^n\le\delta
\), which is equivalent to
\(
n\ge\log\delta/\log(1-pq)
\); taking the least integer proves~\eqref{eq:geometric-quantile}.

Let \(h(x)=-\log(1-x)\).  Convexity of \(h\), together with \(h(0)=0\), gives
\(h(qp)\le qh(p)\), proving the lower residual bound.  Also,
\[
  h(x)=\int_0^x\frac{dt}{1-t}
  \quad\Longrightarrow\quad
  x\le h(x)\le\frac{x}{1-x}.
\]
Therefore
\(
q h(p)/h(qp)\le q[p/(1-p)]/(pq)=1/(1-p)
\), proving the upper residual bound.  Finally, for every \(u>0\),
\(u\le\lceil u\rceil<u+1\); applying \(\log\) with
\(u=\widetilde Q_{p,q}(\delta)\) proves~\eqref{eq:ceiling-residual}.
\end{proof}

\subsection{Proof of Corollary~\ref{cor:retry-crossover}}

\begin{proof}
After \(D\) independent low-cost attempts, failure probability is
\((1-p_\ell)^D\), so reliability is \(1-(1-p_\ell)^D\).  Matching the high-cost
one-attempt reliability requires
\[
  1-(1-p_\ell)^D\ge p_h
  \quad\Longleftrightarrow\quad
  D\ge\frac{\log(1-p_h)}{\log(1-p_\ell)}.
\]
The least integer is~\eqref{eq:retry-crossover}; multiplying by
\(\kappa_\ell\) gives the fixed-block cost condition.  With stopping, attempt
\(d\) is incurred only if the first \(d-1\) attempts fail, which has probability
\((1-p_\ell)^{d-1}\).  Therefore
\[
  \E[\text{cost through depth }D]
  =\kappa_\ell\sum_{d=1}^{D}(1-p_\ell)^{d-1}
  =\kappa_\ell\frac{1-(1-p_\ell)^D}{p_\ell}.
\]
\end{proof}

\end{document}
